% !TeX root = 程序使用手册.tex
%=============================================================================
% 93_changelog.tex -- 更新记录（\input 自 docs/程序使用手册.tex 的
%                     \section{更新记录} 之内）
% 维护规则（.trae/rules/project_rules.md）：每新增/修改一个功能，必须在本表
% 顶部追加一条（日期 / 节点（版本）/ 功能 / 涉及源码 / 验证），并与
% memory-bank/worklog.md 的对应条目一致。
%=============================================================================

本表按时间倒序排列，字段含义与维护要求如下：

\begin{itemize}\itemsep2pt
  \item \textbf{日期}：完成日（与 \texttt{memory-bank/worklog.md} 一致）。
  \item \textbf{节点}：开发阶段号（\texttt{plan.md} 的阶段划分），可视为版本标签。
  \item \textbf{功能}：一句话说明改了什么、为什么改。
  \item \textbf{涉及源码}：主要改动的源文件（不含文档）。
  \item \textbf{验证}：{\bfseries 必须写可复现的验证方式与关键数值}，不能只写``已测试''。
\end{itemize}

{\footnotesize
\begin{longtable}{@{}>{\raggedright\arraybackslash}p{0.075\textwidth}>{\raggedright\arraybackslash}p{0.055\textwidth}>{\raggedright\arraybackslash}p{0.28\textwidth}>{\raggedright\arraybackslash}p{0.24\textwidth}>{\raggedright\arraybackslash}p{0.29\textwidth}@{}}
\caption{MixNSSolver 更新记录}\label{tab:changelog}\\
\toprule
日期 & 节点 & 功能 & 涉及源码 & 验证 \\
\midrule
\endfirsthead
\multicolumn{5}{@{}l}{\small\itshape 表~\ref{tab:changelog}（续）}\\
\toprule
日期 & 节点 & 功能 & 涉及源码 & 验证 \\
\midrule
\endhead
\bottomrule
\endlastfoot
2026-10-09 & 14 &
\textbf{Zhang 2011 界面封闭（阶段 14，全部 opt-in）}：把流固界面封闭改为按文献实现。
① 界面\textbf{速度} \texttt{iface\_velocity=zhang}：两侧共用同一界面速度（Eq.22），法向由
两侧法向应力平衡给出（Eq.25，倒数距离加权调和平均），切向由 Ochoa-Tapia--Whitaker
应力跳变给出（Eq.26，含过量黏性 $\beta$ 与惯性 $\beta_1$ 项，二次式取稳定根）——这正是
此前 ``必须用 balance'' 的文献替代（$\beta\to0$ 退化为应力连续的调和平均，$\beta\to\infty$
为无滑移）。② 界面\textbf{温度} \texttt{iface\_t\_model=zhang}：Eq.27--29，流体侧看到的
界面温度是孔隙率\textbf{体积平均}
$T_{fl}=\langle T\rangle^p=\varepsilon\langle T_f\rangle^f+(1-\varepsilon)\langle T_s\rangle^s$
（既非 $T_f$ 也非 $T_s$ 单相温度），热流按面积比（孔隙率）分给两相（LTNE；LTE 下退化为原
Dirichlet）。③ 新增开关 \texttt{iface\_p\_model}(grad0/dirichlet/momentum)、
\texttt{iface\_slip}(off/bj/noslip)、\texttt{iface\_beta}/\texttt{beta1}/\texttt{df\_ratio}，
并把闭式解集中到 \texttt{mod\_iface\_law}。\textbf{默认全关}（\texttt{iface\_velocity=full}、
\texttt{iface\_t\_model=off}），旧路径位级不变。 &
\texttt{src/common/mod\_iface\_law.f90},
\texttt{src/common/mod\_reference\_state.f90},
\texttt{src/unstructured/mod\_uns\_bc.f90},
\texttt{src/unstructured/mod\_uns\_simple.f90},
\texttt{src/main.f90} &
\texttt{make iface\_law\_test}：$7\to11$ 项全 PASS（新增 Eq.25 调和平均与
$d_f\!\to\!\infty/d_p\!\to\!0$ 极限、Eq.26 的 $\beta$/$ \beta_1$/无实根退化、Eq.22 矢量重构、
Eq.27--29 的 $\varepsilon T_{fi}+(1-\varepsilon)T_{si}=T_{fl}$ 与能量守恒
$k_{fe}\Delta T_f+k_{se}\Delta T_s=k_f\Delta T_{fl}$、LTE 退化为串阻调和平均）；
\texttt{C\_P\_test} \texttt{iface\_velocity=zhang}（$d_f/d_p=50$、$\omega=0.4$）30 耦合步
无发散：$|u|_{\max}$ $0.40\to5.4$ m/s（物理 BJ 滑移量级，对比 \texttt{balance} 基线
$0.30$ m/s、\texttt{iface\_slip=bj} 的 $42$ m/s 发散），界面质量配平
$\dot m$ 残差 $6.9\times10^{-4}$（通量的 $0.3\%$），结构侧界面 $p=1.0148\times10^5$ Pa
与基线一致；\textbf{回归}：旧配置重跑（\texttt{mix.control}，30 耦合步）与基线日志的
全部数值诊断行 \texttt{diff} \textbf{为 0}（仅多一行 \texttt{iface\_t\_model} 报告、MPI
输出交错顺序不同），即旧路径位级不变；60 耦合步续跑（\texttt{mix\_f\_zhang60}）滑移
$0.40\to9.3$ m/s \textbf{单调、增速递减}（$+0.167$ 步$^{-1}$$\to$$+0.126$ 步$^{-1}$）
⇒ 稳定但收敛慢（$d_f$ 进入交换协议后可望完全收敛）；串行/MPI 全构建 rc=0。 \\
\textbf{修复非结构能量方程 ``定温 $+$ 出流'' 边界面的双计焓假热汇}：边界循环对
Dirichlet 面无条件写 \texttt{rhs += DT*T\_face - FT*T\_face}，而出流
（$FT>0$）时迎风值本就是内部 $T_{c0}$（已由 \texttt{ap += max(FT,0)} 隐式计入），
多出的 $-FT\,T_{\rm face}$ 是与流出焓同量级的\textbf{假热汇}，把贴界面第一层单元压到
所有边界值之下（\texttt{C\_P\_test}：入口/界面/绝热壁全 300 K，板块仍出现
$T_{\min}=286$ K）。改为出流项只在入流（$FT<0$）保留：
\texttt{rhs += DT*T\_face - min(FT,0)*T\_face}。 &
\texttt{src/unstructured/mod\_uns\_simple.f90}\newline
(\texttt{temperature\_assembly}) &
\textbf{① 独立复现判据}：把 \texttt{C\_P\_test} 的耦合界面换成普通
\texttt{velocity-inlet} 出流（同网格/同流量/同 300 K，绕开全部耦合代码）后仍逐位复现
\texttt{it=1 dT\_max=7.21}、\texttt{it=50 Tmin=286.0}（耦合 run6 为 7.18/286.1）
$\Rightarrow$ 与界面交换层无关。\\
\textbf{② 修复后同一算例}：\texttt{it=1 dT\_max=2.84e-2}，
$T_{\min}=T_{\max}=300.0$ K（均匀 300 K 初场 $+$ 全 300 K 边界的不变量恢复）。\\
\textbf{③ 耦合 \texttt{C\_P\_test} run7}（30 步，\texttt{du\_max}$\to1.4\times10^{-5}$）：
两侧界面温度 $300.1/286.5\to$ \textbf{300.13/300.12} K，板内 $T\in[299.6,303.5]$ K
（界面高于 300 K 的温升即发汗冷却吸热，主流 309 K）；速度/压力场与修复前
\textbf{逐位相同}（$\overline{|u|}=0.1750$、$|u|_{\max}=0.3154$、表压 115.5--174.9 Pa）
$\Rightarrow$ 只去掉假热汇、不动量解。\\
\textbf{④ 回归}：\texttt{cases/porous\_plug} darcy 与提交日志逐位一致（101 步、
\texttt{mass-imbal}=1.082e-18、$|u|_{\max}=1.03566\times10^{-1}$）；
\texttt{regress/m6wing} PASS。\\
\textbf{遗留}：动量方程同址（\texttt{mod\_uns\_simple.f90:1058}）与
\texttt{BC\_FARFIELD} 入流支有同形项 $-F u_f$；本算例该假源项
$\dot m|u_t|\approx4.3\times10^{-4}$ N 仅为 Darcy 阻力 $3.8\times10^{-2}$ N 的
$\sim1\%$，本次未动，待专案评估。 \\
2026-10-09 & 11 &
\textbf{对拍采样口径修正}（不改源码）：\texttt{cases/betchen} 的 PLUG 对比脚本原用
\texttt{np.interp} 把数字化参考 CSV 中落在我们第一个单元内部（$x/H<0.075$，如
0.032/0.064）的采样点\textbf{截断}成第一列单元值，而 VTU 只有单元中心数据，
于是``入口 u''被读成单元值（1.483）而非 BC 施加的\textbf{面值}（精确 1.5\,$U_0$）；
同时 \texttt{images/plug\_compare\_ref.png} 一直停在旧（均匀入口）VTU 上未重绘。
现 \texttt{vs\_ref.py}/\texttt{compare\_plug.py} 剔除这些点并打印 \texttt{note}，
\texttt{inlet\_profile.py} 增印同样采样的解析离散积分（NY=21 为 0.99449，而非 1.0），
两个绘图脚本一律由当前 VTU 重新生成。另用 x 方向加密 3 倍
（\texttt{gen\_plug.py <out>.cas x 60,40,60}，$\Delta x$ 0.15H$\to$0.05H）
证明第一列偏差 $-1.2\%\to-0.2\%$ 是入口单元 $O(\Delta x)$ 误差而非 BC 误差。 &
\texttt{cases/betchen/README.md} \S3.4.1（新）,
\texttt{cases/betchen/\{vs\_ref,compare\_plug,inlet\_profile\}.py} &
修正后的同一二进制 UNI$\to$PAR 对拍：入口列 $\max|u-6\xi(1-\xi)|$
\textbf{0.398$\to$0.018}（平顶中心 1.102 / 壁面格 0.328 $\to$ 抛物 1.483 / 0.136）；
整体中心线 u L2 0.579\%$\to$\textbf{0.494\%}（Da=1e-2）、2.266\%$\to$2.271\%
（Da=1e-3，由界面过渡层主导），p L2 2.205\%$\to$\textbf{2.157\%} / 1.790\%$\to$\textbf{1.741\%}；
加密网格第一列中心 $u/U_0$ 1.4826$\to$1.4965、$x/H{=}0.075$ 列 1.5002
（$\max|u-6\xi(1-\xi)|$=0.0017），网格加密与入口单元误差按一阶收敛
\\
2026-10-08 & 11 &
新增 \texttt{velocity-inlet-parabolic} 边界（\texttt{BC\_VINLET\_PARAB}=10）：沿面内法向
直接施加充分发展抛物剖面 $|u_f|=6U_{mean}\xi(1-\xi)$、$\xi=(x_{axis}-origin)/span$，
剖面均值恰为 $U_{mean}$（流量与均匀入口一致，只改形状），供论文型``充分发展入口''算例
使用；动量/温度/LTNE 固相识别/BC 报告全部接入。\textbf{动机}：Betchen PLUG 算例入口
原用均匀速度，第一列平顶 ($u/U_0\approx1.10$、壁面格 0.33) 使 PLUG
Da=1e-3 中心线 u L2=7.0\%（入口采样点偏差 1.10 vs 1.50）。 &
\texttt{src/unstructured/mod\_uns\_control.f90},
\texttt{src/unstructured/mod\_uns\_bc.f90},
\texttt{src/unstructured/mod\_uns\_simple.f90},
\texttt{cases/betchen/\{plug\_dae2,plug\_dae3,plug\_hir\}.control} &
\texttt{cases/betchen/README.md} \S3.4/\S4.4：
同一二进制只改入口一行 ⇒ 入口列 $\max|u-6\xi(1-\xi)|$ \textbf{0.398→0.018}、
中心线 u L2 0.579\%→\textbf{0.494\%}（Da=1e-2；max$|\Delta u|$ 0.0150→0.0135，
位置在多孔界面过渡层 $x/H$=3.05）、2.266\%→2.271\%（Da=1e-3；max$|\Delta u|$=0.097
在 $x/H\approx$2.96）、p L2 2.205\%→\textbf{2.157\%} / 1.790\%→\textbf{1.741\%}；
入流剖面第 4 列 ($x/H$=0.525) 与解析 $6\xi(1-\xi)$ 逐点差 $\le0.004$；
\textbf{采样口径修正见 2026-10-09 条}（旧稿的 8.06\%→0.63\% 系参考采样点被
\texttt{np.interp} 截断到入口单元所致）；均匀入口路径\textbf{逐位无回归}
（新二进制+旧入口行重跑 \texttt{plug\_dae3} 的 VTU md5 与旧产物完全一致），
MPI np2 vs 串行 max$|\Delta u|$=2.8e-9
\\
2026-10-08 & 11 &
修复流体/多孔界面压力面值的不一致与 Darcy 汇系数：新增 \texttt{kink\_face\_pressure}
（界面两侧各自单侧二次重构取平均，仅当两侧单元类型不同时替换距离加权插值），
界面压力对"C0 连续、斜率跳变"的压力场恢复一致；\texttt{momentum\_assembly} 的
Darcy 汇由 $(\mu/\varepsilon)/K\,\mathbf u$ 改为 $\mu/K\,\mathbf u$（体积平均动量方程
中 $1/\varepsilon$ 只作用于时间/对流/粘性项，Darcy 汇用表观速度） &
\texttt{src/unstructured/mod\_uns\_fields.f90},
\texttt{src/unstructured/mod\_uns\_simple.f90} &
\texttt{cases/betchen/abtest\_interface\_pressure.sh}（\texttt{git stash} 出 pre-fix
二进制，7 个算例 PRE/POST 对拍，逐列/全场比较）：PLUG Da=$10^{-2}$ 与数字化参考解
$u$ 的 L2 由 $2.94\%\to0.58\%$、$p$ 的 L2 $15.8\%\to2.6\%$；Da=$10^{-3}$ 由
$13.7\%\to2.3\%$、$21.8\%\to1.8\%$；界面逐列抖动 $0.147\to0.018$、
$0.781\to0.074$；Re$_H=1000$ 抖动 $0.142\to0.022$（$\pm10\%$ 锯齿消失）；
\texttt{fluid}（同网格无多孔块）两版 VTU 与求解日志 md5 完全相同、全字段差分 0
（no-op）；\texttt{porous\_plug} 一维塞 dp/dx $-461.5\to-184.6$ Pa/m
（解析 $-184.6$，误差 $150.0\%\to0.0\%$） \\
2026-10-08 & 文档 &
建立本手册（\texttt{程序使用手册.tex}）、附录 A 参数总表、
\texttt{control.ec.template}、本更新记录；补 \texttt{cases/couple\_channel/README.md}
（布局/量纲匹配/缺陷/修法/结果/复现）；修 \texttt{compare\_iface.py} 取窗宽度缺陷 &
（仅文档与工具脚本） &
\texttt{cd docs \&\& xelatex 程序使用手册.tex} 两遍无 Error；模板用与
\texttt{read\_parameter\_ec} 同构的 namelist 读入测试通过；
\texttt{python3 compare\_iface.py} 复现界面压差 4.1 Pa \\
2026-10-07 & 11 &
界面类型自动分派表：\texttt{Interface\_FACE\_TYPE} 增 \texttt{cz\_type/loc\_face/iface\_type}，
新增 \texttt{classify\_interface} / \texttt{iface\_exchange\_recipe} /
\texttt{report\_interface\_dispatch}；纯分类与报告，不改数值行为 &
\texttt{src/common/mod\_interface.f90},
\texttt{src/coupling/mod\_interface\_match.f90},
\texttt{src/coupling/test\_match.f90},
\texttt{src/structured/mod\_struct\_grid.f90},
\texttt{src/unstructured/mod\_uns\_driver.f90} &
\texttt{mpirun -np 1 bin/match\_test}（\texttt{grid\_BC} 拷贝）250/250 全
\texttt{comp-fluid<->lowspeed-fluid}、\texttt{unclassified=0}、断言 PASS；
\texttt{couple\_channel} np2 判为 fluid、\texttt{couple\_porous} np2 判为 porous；
\texttt{regress/m6wing/run\_regression.sh} 全 IDENTICAL（\texttt{flow3d.dat} md5
\texttt{dc134a2d196422043ecad7c86ac8f898}） \\
2026-10-07 & 11 &
修复 uns 绝对压力约 $-250$ Pa 偏置：根因是 \texttt{momentum\_assembly} 内部面对流系数
\texttt{F = ctrl\%rho*fld\%flux(i)} 把已是质量通量（kg/s）的 \texttt{fld\%flux}
重复乘 $\rho$；并给 \texttt{pressure-outlet} 加出口压力爬升 &
\texttt{src/unstructured/mod\_uns\_simple.f90},
\texttt{src/unstructured/mod\_uns\_bc.f90} &
\texttt{cases/couple\_porous/README.md} \S5 四组对照（$\rho$ 倍放大消失）；
\texttt{regress/m6wing/run\_regression.sh} 位级 IDENTICAL \\
2026-10-07 & 11 &
\texttt{pval} 取值敏感性实测后纠正文档：常密度下 \texttt{pval} 只是表压锚点，
但正负不对称（正侧约 $+15$ Pa、负值到约 $-200$ Pa，$+100$ 及以上必发散） &
（仅文档：\texttt{cases/couple\_porous/README.md}） &
各值分别重跑 5 次统计收敛/发散；慢化 \texttt{inlet\_ramp}（300/1000/3000）对
$+100$ 无效 \\
2026-10-06 & 11 &
交换层改为 {\bfseries Dirichlet--Neumann 分区特征界面}：结构侧只接收背压
$p_b$ 并用线性化 Riemann 反射构造 ghost，非结构侧只施加速度 Dirichlet 且界面压力零梯度；
删除压力重锚定；背压松弛 $\alpha = 0.3\min(1, iter/\texttt{iface\_ramp})$ &
\texttt{src/main.f90},
\texttt{src/coupling/mod\_coupling\_exchange.f90},
\texttt{src/coupling/mod\_interface\_exchange.f90} &
\texttt{cases/couple\_channel} np2、\texttt{n\_couple}=1000：无
\texttt{OverLimit}、约 100 迭代后冻结；界面压差 4.1 Pa（旧方案为 kPa 级冻结跳变）；
结构面/内部与 uns 下游 $u_x$ 逐位相同 \\
2026-10-06 & 11 &
C2 可压缩--多孔跨组界面验证并修两处缺陷：多 cell zone 的 cell-id 排序
（分割方向索引最慢）与取窗宽度（$\le$ 一个单元、用绝对窗） &
\texttt{cases/couple\_porous/}（工具与 README）、
\texttt{src/unstructured/mod\_uns\_cas\_reader.f90} &
\texttt{check\_bed\_gradient.py} 带回归守卫（zone 2 不是 $x<100$ mm 半区即报错）；
界面跳变由``幻影 112 Pa''回到物理值 \\
2026-10-06 & 11 &
彻底删除 \texttt{IF\_InnerFlow} / \texttt{IF\_TurboMachinary} 全套机制（结构侧不再
自带内流/叶轮机械分支） &
\texttt{src/structured/mod\_struct\_*},
\texttt{src/main.f90} &
\texttt{regress/m6wing/run\_regression.sh} 位级 IDENTICAL \\
2026-10-06 & 10 &
流场自动保存 + 耦合联合重启：\texttt{save\_interval}（每 N 个耦合步与两侧同步存档）、
\texttt{couple\_restart}（结构侧强制从 \texttt{flow3d.dat} 续算，非结构侧读耦合状态） &
\texttt{src/main.f90},
\texttt{src/structured/mod\_struct\_driver.f90} &
\texttt{couple\_porous}/\texttt{couple\_channel} np2 跑满 \texttt{n\_couple} 且
\texttt{save\_interval} $\le$ \texttt{n\_couple} 时 \texttt{flow3d.dat} 落盘；
二次启动从存档续算，界面量连续 \\
2026-10-06 & 9 &
多孔介质热学收尾：LTNE（流体/固体双温度）＋热弥散（\texttt{disp\_l}/\texttt{disp\_t}）、
对角各向异性渗透率（\texttt{perm\_xx/yy/zz}）、Beavers--Joseph 界面滑移
（\texttt{bj\_alpha}）、VC 标记自动识别块类型 &
\texttt{src/unstructured/mod\_uns\_simple.f90},
\texttt{src/unstructured/mod\_uns\_bc.f90},
\texttt{src/unstructured/mod\_uns\_control.f90} &
\texttt{cases/ltne}（1D 发汗冷却解析解）、\texttt{cases/graetz2d}（2D Graetz）、
\texttt{cases/porous\_plug}（1D 压降）解析/半解析对比；低速-多孔界面确认复用 BJ Robin 条件 \\
2026-10-05 & 5--6 &
弱耦合框架打通：\texttt{mix.control} 参考状态、\texttt{MPI\_Comm\_split} 分组、
\texttt{main.f90} 弱耦合循环、接口量交换先转 SI；
新增 \texttt{outflow}（充分发展出口）与 \texttt{mass-flow-inlet} 起动 ramp &
\texttt{src/main.f90},
\texttt{src/common/mod\_reference\_state.f90},
\texttt{src/common/mod\_interface\_units.f90},
\texttt{src/unstructured/mod\_uns\_control.f90} &
常量场 MPI 端到端测试（两组交换值逐位一致）；
\texttt{cases/channel} 槽道起动稳定（\texttt{inlet\_ramp}） \\
2026-10-04 & 8 / 10 &
瞬态 PISO（\texttt{transient}/\texttt{dt}/\texttt{n\_correct}/\texttt{time\_scheme}）
与 PIMPLE 稳定化（\texttt{pimple}/\texttt{n\_outer\_iter}）；能量方程（被动标量）
＋ Boussinesq 浮力（\texttt{cp}/\texttt{k\_cond}/\texttt{beta}/\texttt{gravity}/
\texttt{boussinesq}）＋ 多孔 Darcy--Forchheimer（LTE/LTNE） &
\texttt{src/unstructured/mod\_uns\_simple.f90},
\texttt{src/unstructured/mod\_uns\_control.f90},
\texttt{src/unstructured/mod\_uns\_output.f90} &
\texttt{cases/cavity\_transient}（PISO，$\max|u|$ 与稳态一致量级）、
\texttt{cases/cavity\_temp}（被动标量：$\max|u|$ 不变、$0\le T\le1$）、
\texttt{cases/cavity\_boussinesq}、
\texttt{cases/natconv}（多孔自然对流）、\texttt{cases/rayleigh\_benard}（2026-10-08） \\
2026-10-04 & 4 &
交界面几何匹配：最近点/投影 ＋ 双线性插值权重，双向映射（struct 面 $\leftrightarrow$
uns 面） &
\texttt{src/coupling/mod\_interface\_match.f90},
\texttt{src/common/mod\_interface.f90} &
\texttt{bin/match\_test} 逐面比对（250/250 命中、权重和为 1） \\
2026-10-03 & 3 &
非结构求解器迁移进树；Fluent \texttt{.cas} 的 interface 面 zone 与 cell zone 解析、
\texttt{.control} 增 \texttt{cell\_zone}、新增 \texttt{mass-flow-inlet} 与
\texttt{pressure-far-field}，修复 far-field 零解缺陷 &
\texttt{src/unstructured/*} &
\texttt{regress} 与源求解器位级一致；\texttt{grid\_BC/cavity} 算例回归 \\
2026-10-03 & 2 &
结构求解器原样搬迁 + 模块化封装（\texttt{mod\_struct\_*} 前缀、全局模块名唯一、
共享 \texttt{mod\_precision}/\texttt{mod\_constants}）；新增 \texttt{BC\_Interface} &
\texttt{src/structured/*},
\texttt{src/common/*} &
\texttt{regress/m6wing/run\_regression.sh} 每批封装后位级 IDENTICAL \\
2026-10-03 & 1 &
顶层 Makefile：\texttt{ser}/\texttt{mpi} 双构建树、\texttt{.mod} 集中目录、
自动依赖（\texttt{-MMD -MP}）、\texttt{DEBUG=1} &
\texttt{Makefile} &
\texttt{make all mpi} 全绿；\texttt{make help} 列出全部目标 \\
\end{longtable}
}
